\begin{table}[htbp]
\centering
\caption{Data overview for the three cell types analyzed.}
\label{tab:supp-data}
\begingroup
\small
\singlespacing
\setlength{\tabcolsep}{4pt}
\renewcommand{\arraystretch}{1.3}
\begin{tabular}{@{}>{\raggedright\arraybackslash}p{3.3in}rrr@{}}
\toprule
 & Bn & CD4 T cell & CD8 T cell \\
\midrule
\rowcolor{gray!14}\multicolumn{4}{@{}l}{\textbf{Single-cell data used for the pseudobulk correlation (after quality control)}} \\
Individuals (passing RNA and ATAC cell-count filters) & 335 & 372 & 346 \\
scRNA-seq cells & 160,577 & 916,643 & 460,439 \\
scATAC-seq cells & 123,580 & 360,343 & 210,077 \\
scRNA-seq cells per individual, median (IQR) & 436 (315--619) & 2,336 (1,655--3,150) & 1,200 (763--1,760) \\
scATAC-seq cells per individual, median (IQR) & 348 (236--464) & 898 (618--1,278) & 510 (320--845) \\
Genes in the pseudobulk RNA matrix & 36,326 & 36,326 & 36,326 \\
Peaks in the pseudobulk ATAC matrix & 338,036 & 338,036 & 338,036 \\
\rowcolor{gray!14}\multicolumn{4}{@{}l}{\textbf{cis-QTL summary statistics}} \\
eQTL sample size & 405 & 406 & 405 \\
caQTL sample size & 385 & 385 & 384 \\
Peaks with at least one caQTL at \(P<10^{-8}\) & 3,449 & 6,110 & 3,847 \\
\rowcolor{gray!14}\multicolumn{4}{@{}l}{\textbf{Candidate peak--gene pairs and the Mendelian randomization analysis set}} \\
Candidate gene--peak pairs (All\_GenePeak, TSS--peak distance \(<1\)~Mb) & 3,814,826 & 5,146,441 & 4,505,251 \\
With caQTL and eQTL variants in common (at least 3 shared variants) & 46,190 & 109,961 & 62,136 \\
Gene and peak present in both pseudobulk and WNN data & 40,389 & 86,204 & 51,649 \\
After removing pairs with a missing p-value in any base method & 40,153 & 85,782 & 51,123 \\
Final analysis set (IVW residual variance \(\le\) 10) & 40,087 & 85,527 & 50,931 \\
Genes / peaks in the final analysis set & 8,149 / 2,951 & 10,701 / 5,338 & 9,427 / 3,327 \\
Variants (instruments) per pair in the IVW fit, median (IQR) & 22 (10--40) & 22 (10--43) & 22 (9--41) \\
\rowcolor{gray!14}\multicolumn{4}{@{}l}{\textbf{Samples for the marginal correlation methods}} \\
Pseudobulk correlation: individuals & 335 & 372 & 346 \\
WNN correlation: metacells (external PBMC multiome) & 80 & 80 & 80 \\
\bottomrule
\end{tabular}
\par\medskip
\begin{minipage}{\linewidth}
\footnotesize
Bn, naive B cells (CIMA L4 type Bn\_TCL1A); CD4 T cell, naive CD4\(^+\) T cells (CD4\_Tn\_CCR7);
CD8 T cell, naive CD8\(^+\) T cells (CD8\_Tn\_CCR7). The Chinese Immune Multi-Omics Atlas (CIMA;
Yin et al., Science 2026; OMIX009786) comprises scRNA-seq and scATAC-seq data from
428 Chinese adults annotated into 73 L4 immune cell types. The cis-QTL were mapped within
1~Mb of the gene transcription start site or the peak midpoint, using genotypes of the
whole cohort (5,051,397 variants with minor allele frequency \(>0.1\)); the numbers
in the table are specific to the three cell types analyzed here.
Three sample definitions are used: the QTL sample sizes (summary statistics used by the
Mendelian randomization methods), the individuals in the pseudobulk matrices (those passing the
RNA and ATAC cell-count filters; pseudobulk correlation), and 80 metacells of an external PBMC
multiome data set (WNN correlation). The caQTL and eQTL data come from the same cohort, so
the two QTL data sets overlap in samples.
Peaks with at least one caQTL at \(P<10^{-8}\) are counted from the significant caQTL
files used to select instruments; this is a nominal threshold, not the false-discovery-rate-defined
caPeak count of CIMA.
Candidate pairs are All\_GenePeak pairs with at least 3 variants shared between the clumped
caQTL and the eQTL data, whose gene and peak are present in both the pseudobulk and WNN data.
Pairs with a missing p-value in any base method (IVW, MR-Egger, SMR, WNN, pseudobulk) or with
an IVW residual variance above 10 were removed. IQR, interquartile range.
\end{minipage}
\endgroup
\end{table}
